\documentclass[11pt]{article}
\usepackage[margin=2.7cm]{geometry}
\usepackage{amsmath,amssymb}
\usepackage[hidelinks]{hyperref}

\newcommand{\Op}{\operatorname{Op}^{w}}
\newcommand{\pd}{\partial}

\title{A Gentle Introduction to Weyl Symbols}
\author{}
\date{}

\begin{document}
\maketitle

\section{The main idea}

A differential operator acts on a function. A \emph{symbol} is a way to
describe that operator using an ordinary function of position and
wavenumber. It is useful because many questions about an operator---what
waves it supports, how it propagates them, and how two operators combine---
become easier to see in this position--wavenumber picture.

For a wave $u(x)=e^{i\xi x}$, the derivative measures its wavenumber:
\begin{equation}
  D u=\xi u,\qquad D=-i\pd_x.
\end{equation}
So $D$ is represented by the variable $\xi$. For example, the constant
coefficient operator $D^2+k^2$ has symbol $\xi^2+k^2$. Its symbol immediately
shows which wavenumbers make $(D^2+k^2)u=0$.

When coefficients vary with position, position and wavenumber both matter.
The operator $a(x)D$ has a symbol involving both $a(x)$ and $\xi$. There is a
small complication: multiplication by $a(x)$ and differentiation do not
commute. Different choices for their ordering give slightly different
symbols. Weyl's choice treats position and wavenumber symmetrically.

\section{Weyl quantization}

Given a function $p(x,\xi)$, its Weyl operator is
\begin{equation}\label{eq:weyl}
  (\Op(p)u)(x)=\frac{1}{2\pi}\iint
  e^{i(x-y)\xi}p\!\left(\frac{x+y}{2},\xi\right)u(y)\,dy\,d\xi.
\end{equation}
The key feature is the midpoint $(x+y)/2$: the coefficient is sampled
halfway between where the input and output are evaluated. This makes Weyl
quantization symmetric in position and wavenumber.

Some basic examples are
\begin{align}
  \Op(a(x))u &=a(x)u,\\
  \Op(\xi)u &=Du=-i u_x,\\
  \Op(a(x)\xi)u &=\frac12\bigl(aD+Da\bigr)u
    =-i\left(a u_x+\frac12a_xu\right).
\end{align}
Thus the symbol $a(x)\xi$ represents the \emph{symmetrically ordered}
operator, not simply $aD$. In fact,
\begin{equation}
  aD=\Op\!\left(a(x)\xi+\frac{i}{2}a_x(x)\right).
\end{equation}
The extra term records the effect of the ordering. For real $p$, the Weyl
operator is symmetric (subject to the usual domain and boundary conditions).

\section{Combining symbols}

Applying one operator after another corresponds to combining their symbols
with a product called the \emph{Weyl product}, written $p\#q$. Its first terms
are
\begin{equation}\label{eq:moyal}
  p\#q=pq+\frac{i}{2}
  \left(p_xq_\xi-p_\xi q_x\right)
  -\frac18\left(p_{xx}q_{\xi\xi}-2p_{x\xi}q_{\xi x}
  +p_{\xi\xi}q_{xx}\right)+\cdots.
\end{equation}
The first term is ordinary multiplication. The remaining terms correct for
the fact that position-dependent coefficients and derivatives do not
commute. If the coefficients vary slowly compared with a wavelength, these
corrections are often small, so ordinary multiplication gives a useful first
approximation.

For example, the symbol of a commutator is, to leading order,
\begin{equation}
  \sigma\bigl([\Op(p),\Op(q)]\bigr)
  \simeq i\left(p_xq_\xi-p_\xi q_x\right).
\end{equation}
This expression is the phase-space Poisson bracket (with the convention
shown here). It tells us that operator order matters most where the two
symbols vary across phase space.

\section{A variable-coefficient wave operator}

Consider the one-dimensional vertical operator
\begin{equation}
  B=\pd_x\bigl(g(x)\pd_x\bigr)+V(x),\qquad g(x)>0.
\end{equation}
Since $D=-i\pd_x$, this is $B=-DgD+V$. Its exact Weyl symbol is
\begin{equation}\label{eq:variable}
  b(x,\xi)=-g(x)\xi^2+V(x)-\frac14g_{xx}(x).
\end{equation}
The leading part $-g\xi^2+V$ is the local dispersion relation. The smaller
$-g_{xx}/4$ term is a correction caused by the variation of the coefficient
and by symmetric operator ordering. In a slowly varying medium, the
principal symbol often explains the main wave behavior, while the lower
order terms describe corrections to phase and amplitude.

In terrain-following Helmholtz analysis, the same idea applies with position
$\eta$ and vertical wavenumber $q$. The symbol of the vertical operator gives
the local dispersion relation at each range, while its range variation
couples forward and backward waves. Weyl notation keeps the operator ordering
visible and provides a systematic way to calculate those coupling and
higher-order terms.

\paragraph{Convention.} This note uses $D=-i\pd_x$, the kernel
$e^{i(x-y)\xi}$, and the Poisson bracket
$\{p,q\}=p_xq_\xi-p_\xi q_x$. Changing the Fourier or bracket convention
changes some signs together; the underlying operator statements are
unchanged.

\end{document}
